\newcommand{\C}{\mathbb{C}}
\newcommand{\R}{\mathbb{R}}
\newcommand{\Z}{\mathbb{Z}}
\newcommand{\Tr}{\operatorname{Tr}}
\newcommand{\Impart}{\operatorname{Im}}
\newcommand{\AF}{\operatorname{AF}}
\newcommand{\AX}{\operatorname{AX}}
\newcommand{\GL}{\operatorname{GL}}

\examyear{2012}
\examera{平成24年}
\examfield{解析}
\examgroup{B}
\examnumber{9}
\examtitle{平成24年 B 第9問}
\examtag{複素解析}
\examtag{正則関数}
\examtag{巾級数}

\(i\) を虚数単位とし，複素数係数の巾級数
\[
  f(z)=i+\sum_{n=1}^{\infty}a_n z^n
\]
を考える。\(f(z)\) は単位円板
\[
  D=\{z\in\C\mid |z|<1\}
\]
上正則であり，各 \(z\in D\) に対し \(f(z)\in H=\{w\in\C\mid \Impart w>0\}\) であるとする。
\(f\) の実部を \(u\)，虚部を \(v\) とし，極座標 \(z=re^{i\theta}\ (0\le r<1,\ 0\le \theta\le 2\pi)\) により
\[
  f(re^{i\theta})=u(r,\theta)+iv(r,\theta)
\]
と表す。このとき次の問いに答えよ。

\begin{enumerate}
  \item 各 \(n>0\) と任意の \(0<r<1\) に対し
  \[
    a_n=\frac{i}{\pi r^n}\int_{0}^{2\pi} v(r,\theta)e^{-in\theta}\,d\theta
  \]
  であることを示せ。
  \item すべての \(n=1,2,\ldots\) に対して \(|a_n|\le2\) を示せ。
\end{enumerate}